\subsection{Resultados parciales}
\label{sec:resultados}

Las pruebas disponibles cubren progresivamente la cadena, desde el acondicionamiento hasta el registro de campo. Los barridos eléctricos verificaron continuidad, ganancia y fase de las rutas internas hasta el ADC. Al combinar la transferencia medida con el modelo nominal del SM-24, la figura~\ref{fig:caracterizacion-electrica} muestra la respuesta esperada del conjunto geófono--electrónica. La comparación distingue el resultado medido de la electrónica de la predicción que incorpora al transductor.

La calibración automática corrigió la continua etapa por etapa antes de adquirir. Las señales de la figura~\ref{fig:calibracion-etapas} convergen hacia sus referencias y los códigos obtenidos se conservan por configuración de ganancia. Una vez calibrado el frente, el PSoC digitaliza y traslada las muestras por las rutas DMA de la figura~\ref{fig:hardware-digital}; la comunicación y la interfaz operan después sobre una captura que ya quedó almacenada localmente.

La coordinación inalámbrica se ensayó de manera exploratoria con tres nodos esclavos. Durante 128 órdenes START se observaron en osciloscopio los flancos GPIO generados por los ESP32 al ejecutar la orden. La separación entre esos marcadores se encontró aproximadamente entre 200 y \SI{400}{\micro s}. El resultado demuestra la operación multiesclavo y proporciona una primera magnitud del arranque distribuido; no mide todavía el inicio efectivo de los ADC ni la deriva de sus relojes.

\begin{figure*}[!tb]
\centering
\includegraphics[width=0.64\textwidth,keepaspectratio]{f_contenido_baja_frecuencia.png}
\caption{Densidad espectral y relación señal--ruido de las capturas, agrupadas por distancia. La zona sombreada identifica la banda baja de interés y la línea vertical la frecuencia natural nominal del SM-24.}
\label{fig:contenido-baja}
\end{figure*}

\subsubsection{Contenido espectral de las capturas}

Se utilizaron ventanas de un segundo inmediatamente antes y después del marcador del impacto. A cada segmento se le restó su media y se le aplicó una ventana Hann; el espectro se normalizó por la energía de esa ventana y la tasa de muestreo. Para cada distancia se obtuvo la mediana de los espectros posteriores, $P_{\mathrm{post}}(f)$, y de los previos, $P_{\mathrm{pre}}(f)$. El cociente representado es
\begin{equation}
\mathrm{SNR}_{\mathrm{ap}}(f)=10\log_{10}\!\left(\frac{P_{\mathrm{post}}(f)}{P_{\mathrm{pre}}(f)}\right).
\end{equation}
Se trata de una relación señal--ruido aparente: la ventana posterior también contiene ruido. Un valor próximo a cero indica que el impacto apenas se distingue del nivel previo en esa banda. La duración de un segundo limita el detalle espectral a una escala del orden de \SI{1}{Hz}; el relleno con ceros sólo interpola la representación. Por ello se interpretan bandas y tendencias entre distancias, no picos estrechos aislados.

Para estimar la banda realmente observable se comparó, por impacto, una ventana posterior al marcador con una ventana de ruido previa. La figura~\ref{fig:contenido-baja} reúne 265 impactos en 15 distancias con al menos tres registros utilizables. La mediana de la relación señal--ruido fue 0,5 dB entre 1 y 5 Hz y 2,1 dB entre 5 y 10 Hz; aumentó a 16,7 dB entre 10 y 20 Hz y a 24,4 dB entre 20 y 50 Hz. Ninguna de las 15 distancias superó 10 dB de mediana por debajo de 10 Hz.

El resultado es extremo a extremo y, por ello, combina fuente, acoplamiento, terreno, geófono y electrónica: no permite atribuir la limitación a un único bloque. Se propone aumentar la energía y la repetibilidad de la excitación en baja frecuencia como siguiente intervención, para evaluar cuánto puede aprovecharse la compensación del receptor y si se extiende la banda útil del registro.

\subsubsection{Registro espacial sintetizado}

La campaña desplazó un receptor entre 21 posiciones de 10 a \SI{50}{m}, con un paso de \SI{2}{m}. En cada posición se repitió el impacto, se alinearon las capturas con la referencia del martillo y se obtuvo una traza representativa. El resultado se ordenó por distancia para formar el waterfall de la figura~\ref{fig:waterfall}.

\begin{figure*}[!tb]
\centering
\includegraphics[width=0.67\textwidth,keepaspectratio]{f_waterfall_4s.png}
\caption{Waterfall sintetizado con 21 posiciones secuenciales y un único receptor desplazado. Se muestran los primeros 4 s posteriores al impacto; cada traza está normalizada para visualizar su evolución temporal.}
\label{fig:waterfall}
\end{figure*}

La progresión temporal observada entre posiciones confirma que la plataforma y el flujo de procesamiento permiten construir un registro espacio--tiempo. No corresponde a 21 receptores simultáneos: cada fila procede de una excitación distinta y la normalización impide comparar amplitudes absolutas. La evidencia alcanzada sostiene la afirmación central de esta etapa: se desarrolló una plataforma orientada a MASW, verificada parcialmente mediante ensayos de sus subsistemas y evidencia preliminar de campo.



\section{Trabajo futuro}

La siguiente intervención será mejorar la excitación en la banda baja de interés. Se propone incorporar una fuente más energética y repetible y comparar su efecto mediante los mismos espectros y estimaciones de relación señal--ruido. La comparación deberá conservar la geometría y documentar el acoplamiento, para reconocer qué mejora se obtiene sobre el registro completo sin atribuir a la fuente todas las limitaciones observadas.

La ampliación del arreglo dependerá de los recursos disponibles para incorporar receptores. Se buscará financiación para varios sensores; si no se obtiene, se continuará mediante adquisición secuencial con un receptor desplazado y una referencia de impacto en cada posición.

El análisis de dispersión, la selección modal y los algoritmos de inversión se reservan para la segunda presentación.

% Apartado de validación futura reservado para definir con el tutor, según comentario del autor.
